\documentclass[10pt,a4paper]{amsart}
\usepackage[margin=19mm]{geometry}
\usepackage{amsmath,amssymb,mathtools}
\usepackage[scr=rsfs,cal=euler]{mathalfa}
\usepackage{xfrac}
\usepackage[unicode,hidelinks,bookmarks=false]{hyperref}
\newcommand{\ie}{i.e.\ }
\newcommand{\cf}{cf.\ }
\newcommand{\resp}{respectively\ }
\newcommand{\Subsection}{\subsection}
\providecommand{\nobreakdash}{\nobreak\hbox{-}}
\setlength{\emergencystretch}{2em}
\multlinegap0pt
\usepackage{amssymb}
\usepackage{bm}
\usepackage{xcolor}
\usepackage{tikz}
\usepackage[shortlabels]{enumitem}
\newtheorem{theorem}{Theorem}
\newcommand{\F}{\numberset{F}}
\newcommand{\fq}{\F_q}
\newcommand{\numberset}{\mathbb}
\def\sym{\textup{sym}}
\title{\textbf{Symmetric Tensor Decompositions over Finite Fields}}
\author{Giuseppe Cotardo, Ferdinando Zullo}
\date{Source: arXiv:2605.12295v1 (CC BY 4.0)}
\begin{document}
\maketitle

\noindent\emph{Source and licence.} \textbf{Symmetric Tensor Decompositions over Finite Fields}. Version arXiv:2605.12295v1 (CC BY 4.0), distributed by arXiv under the Creative Commons Attribution 4.0 International licence, http://creativecommons.org/licenses/by/4.0/, as recorded in the arXiv licence metadata for this exact version and shown on the abstract page https://arxiv.org/abs/2605.12295v1. Full licence notice: \texttt{licenses/arxiv-2605.12295-CC-BY-4.0.txt}.\par\smallskip

\noindent\textit{Editorial scope.} Counted unit: the article's theorem thm:m=2 (source0.tex lines 957--971). The article's complete original proof of it is delivered with it (source0.tex lines 972--1071, one \texttt{proof} environment of 100 lines). No mathematics is written, repaired or re-derived: the statement and the proof are the article's own lines, in the article's own notation, and no retained line is substituted or re-flowed.  Retained uncounted as the source-local context the proof needs: lines 624--630; lines 877--917.  Nothing was omitted from any retained span, so the package carries no omission record.  No retained line contains a citation, so the wrapper carries no bibliography and no bibliography entry.  No retained line contains a figure environment, an image-inclusion command or drawing code, so no image is delivered and the package declares no asset dependency.  Editorial formatting only, in addition: an AMS \texttt{amsart} preamble that restates the article's own declarations and macros used by the retained lines, namely the environments \texttt{theorem} on separate counters, together with the standard packages those lines need.  The retained environments print in this document as Theorem 1 in order of appearance, whereas the article prints the same items with the numbering of its own full document.  The complete native source of the article as distributed in the arXiv e-print for this exact version is bundled unchanged as \texttt{sources/200416/source0.tex}.
\begin{theorem}\label{thm:burg}
For every \(q\) and \(m\), one has
\[
\mu_q(m)\geq 2m-1.
\]
Moreover, $\mu_q(m)=\mu_q^{\mathrm{sym}}(m)=2m-1$ if and only if \(q\geq 2m-2\).
\end{theorem}

\begin{theorem}\label{thm:charactfinal}
Let \(R\) be a positive integer. The following are equivalent.
\begin{enumerate}
    \item \(\mu_q^{\mathrm{sym}}(m)\leq R\).
    \item There exist \(\alpha_1,\ldots,\alpha_R\in\mathbb F_{q^m}^{*}\) such that,
    for every \(i\in\{0,\ldots,m-1\}\), the system \(\Sigma_i\) admits a solution
    in \(\mathbb F_q^R\).
\end{enumerate}

Moreover, let
\[
A=
\begin{pmatrix}
\alpha_1^2 & \alpha_2^2 & \cdots & \alpha_R^2\\
\alpha_1^{q+1} & \alpha_2^{q+1} & \cdots & \alpha_R^{q+1}\\
\vdots & \vdots & & \vdots\\
\alpha_1^{q^{m-1}+1} & \alpha_2^{q^{m-1}+1} & \cdots & \alpha_R^{q^{m-1}+1}
\end{pmatrix}.
\]
If there exist \(\alpha_1,\ldots,\alpha_R\in\mathbb F_{q^m}^{*}\) such that, for every
\(i\in\{0,\ldots,m-1\}\), the enlarged system
\[
\Sigma_i^{*}\colon
\begin{pmatrix}
A\\
A^{[1]}\\
\vdots\\
A^{[m-1]}
\end{pmatrix}
X
=
\begin{pmatrix}
b_i\\
b_i^{[1]}\\
\vdots\\
b_i^{[m-1]}
\end{pmatrix}
\]
admits a solution \(X\in\mathbb F_q^R\), where \(A^{[j]}\) denotes the matrix obtained
from \(A\) by applying the~\(q^j\)-Frobenius entrywise, then $\mu_q^{\mathrm{sym}}(m)\leq R$.
\end{theorem}

\begin{theorem}\label{thm:m=2}
An optimal symmetric tensor-rank decomposition of \(M_{q^2}\) is given by
\begin{equation}\label{eq:3traces1,eta,eta2}
\left\{
\operatorname{Tr}_{q^2/q}(x),\,
\eta\operatorname{Tr}_{q^2/q}(\eta x),\,
\eta^2\operatorname{Tr}_{q^2/q}(\eta^2x)
\right\},
\end{equation}
where \(\eta\in\mathbb F_{q^2}^{*}\) satisfies
\[
\eta^{2q}-\eta^{2q-1}-\eta^{q+1}+\eta^{q-1}+\eta-1\neq 0.
\]
In particular, such an element \(\eta\) always exists.
\end{theorem}

\begin{proof}
By Theorem \ref{thm:burg}, we have $ \mu_q^\sym(2) \geq  \mu_q(2)\geq 3$.
Thus, it is enough to construct a symmetric tensor-rank decomposition of
\(M_{q^2}\) with three elements.

For \(m=2\), we apply Theorem \ref{thm:charactfinal} with \(R=3\). Taking
\(\alpha_1=1\), we seek \(\alpha_2,\alpha_3\in\mathbb F_{q^2}^{*}\) such that, for
\(i\in\{0,1\}\), the enlarged system
\[ \Sigma_i^* \colon \begin{pmatrix}
1 & \alpha_2^2 & \alpha_3^2\\
1 & \alpha_2^{2	q} & \alpha_3^{2q}\\
1 & \alpha_2^{q+1} &  \alpha_3^{q+1}
\end{pmatrix}
\begin{pmatrix} x_1 \\ x_2 \\ x_3 \end{pmatrix} = \begin{pmatrix} \xi^i \\ \xi^{iq} \\ 0 \end{pmatrix},
 \]
admits a solution in
\(\mathbb F_q^3\), where $\xi$ is any generator of $\F_{q^2}$ over $\fq$.
In this case, the coefficient matrix of \(\Sigma_i^*\) is
\[
\begin{pmatrix}
1 & \alpha_2^2 & \alpha_3^2\\
1 & \alpha_2^{2q} & \alpha_3^{2q}\\
1 & \alpha_2^{q+1} & \alpha_3^{q+1}
\end{pmatrix}.
\]
Hence it is enough to choose \(\alpha_2,\alpha_3\) so that this matrix is nonsingular, that is, if its determinant in nonzero. We now choose \(\alpha_2=\eta\) and \(\alpha_3=\eta^2\). The determinant of the
coefficient matrix is then
\[
\det
\begin{pmatrix}
1 & \eta^2 & \eta^4\\
1 & \eta^{2q} & \eta^{4q}\\
1 & \eta^{q+1} & \eta^{2(q+1)}
\end{pmatrix}
=
\eta^{q+2}
\left(
\eta^{2q}-\eta^{2q-1}-\eta^{q+1}+\eta^{q-1}+\eta-1
\right).
\]
Thus, if
\[
\eta^{2q}-\eta^{2q-1}-\eta^{q+1}+\eta^{q-1}+\eta-1\neq 0,
\]
then the coefficient matrix is nonsingular. It remains to show that such an element \(\eta\) exists. If \(q>2\), consider
\[
F(T)=T^{q+2}
\left(
T^{2q}-T^{2q-1}-T^{q+1}+T^{q-1}+T-1
\right)\in \mathbb F_q[T].
\]
A nonzero root in $\F_{q^2}$ of $F(T)$ needs to be a root of $T^{2q}-T^{2q-1}-T^{q+1}+T^{q-1}+T-1$, therefore for $q>2$  the polynomial \(F(T)/T^{q+2}\) has degree \(2q<q^2-q\). Since \(\mathbb F_{q^2}\setminus\mathbb F_q\) has~\(q^2-q\) elements, \(F\) cannot
vanish on all of \(\mathbb F_{q^2}\setminus\mathbb F_q\). If \(q=2\), the determinant reduces to
\[
\eta^4(\eta-1),
\]
using \(\eta^4=\eta\) on \(\mathbb F_4\). Thus any
\(\eta\in\mathbb F_4\setminus\mathbb F_2\) works. For each \(i\in\{0,1\}\), the unique solution of \(\Sigma_i^*\) is given by
Cramer's rule. Namely,
\[\left( \frac{\det\begin{pmatrix}
\xi^i & \alpha_2^2 & \alpha_3^2\\
\xi^{iq} & \alpha_2^{2	q} & \alpha_3^{2q}\\
0 & \alpha_2^{q+1} &  \alpha_3^{q+1}
\end{pmatrix}}{\det\begin{pmatrix}
1 & \alpha_2^2 & \alpha_3^2\\
1 & \alpha_2^{2	q} & \alpha_3^{2q}\\
1 & \alpha_2^{q+1} &  \alpha_3^{q+1}
\end{pmatrix}}, \frac{\det\begin{pmatrix}
\alpha_1^2 & \xi^i & \alpha_3^2\\
\alpha_1^{2	q} & \xi^{iq} & \alpha_3^{2q}\\
\alpha_1^{q+1} & 0 &   \alpha_3^{q+1}
\end{pmatrix}}{\det\begin{pmatrix}
1 & \alpha_2^2 & \alpha_3^2\\
1 & \alpha_2^{2	q} & \alpha_3^{2q}\\
1 & \alpha_2^{q+1} &  \alpha_3^{q+1}
\end{pmatrix}},
\frac{\det\begin{pmatrix}
\alpha_1^2 & \alpha_2^2 & \xi^i \\
\alpha_1^{2	q}  & \alpha_2^{2q}& \xi^{iq}\\
\alpha_1^{q+1}  &   \alpha_2^{q+1} & 0
\end{pmatrix}}{\det\begin{pmatrix}
1 & \alpha_2^2 & \alpha_3^2\\
1 & \alpha_2^{2	q} & \alpha_3^{2q}\\
1 & \alpha_2^{q+1} &  \alpha_3^{q+1}
\end{pmatrix}}
\right).\]
Moreover, \(X_i\in\mathbb F_q^3\). Indeed, the system \(\Sigma_i^*\) is stable under
the \(q\)-Frobenius map, and the solution is unique because \(D\neq 0\). Hence the
Frobenius conjugate of \(X_i\) is again a solution, and therefore it must coincide
with \(X_i\). Thus \(X_i\) is fixed by Frobenius and lies in \(\mathbb F_q^3\). By Theorem~\ref{thm:charactfinal}, the set
\[
\left\{
\operatorname{Tr}_{q^2/q}(x),\,
\eta\operatorname{Tr}_{q^2/q}(\eta x),\,
\eta^2\operatorname{Tr}_{q^2/q}(\eta^2x)
\right\}
\]
gives a symmetric tensor-rank decomposition of \(M_{q^2}\). This and the lower
bound~\(\mu_q^{\sym}(2)\geq 3\), prove the claim.
\end{proof}

\end{document}
